\section{结论}
% PhyPrompt employs a concise two-stage pipeline: first, supervised fine-tuning on a physics-focused Chain-of-Thought corpus, then reinforcement learning with a dynamic reward curriculum that gradually shifts emphasis from semantic fidelity to commonsense physics. This process automatically refines text-to-video prompts for substantially improved physical realism achieving a 40.8\% combined semantic-and-physical success rate, while generalizing zero-shot across diverse generators.
我们提出 PhyPrompt，这是一个基于强化学习的框架，可自动增强文本到视频提示词，从而生成更符合物理规律的视频。通过两阶段训练，即先进行以物理为中心的思维链监督微调，再使用带动态奖励课程的强化学习，PhyPrompt 在保持语义保真度的同时显著提升了物理真实性。
核心实证结果表明，只要通过合适的课程进行优化，语义遵循度和物理常识并不必然相互竞争。PhyPrompt-7B 在 VideoPhy2 基准上达到 47.8\% 的语义遵循度和 66.8\% 的物理常识，联合成功率为 40.8\%。关键在于，物理常识从 55.8\% 提升到 66.8\%，并没有牺牲语义质量；相反，语义遵循度也从 43.4\% 提升到 47.8\%，使联合表现相较基线提示词提高 8.6 个百分点。
这种协同改进来自动态课程机制，它能够发现单目标优化难以到达的提示词空间区域。课程在训练早期建立语义支架，再通过物理约束进行细化，使两个目标相互强化而不是相互抵消。消融研究证实了这一点：动态课程不仅优于静态平衡加权，也在各自指标上优于专门的单目标训练。除性能提升外，PhyPrompt 还展示了跨不同生成器架构（Lavie、VideoCrafter2、CogVideoX）的真正零样本迁移，无需为每个生成器单独微调即可带来一致改进。这种与模型无关的设计使 PhyPrompt 成为物理感知视频生成中实用且可部署的方案。
我们的结果也为生成式 AI 中的多目标优化提供了更广泛的启示：与单纯扩大参数规模相比，带有直接任务反馈的领域专用训练、精心设计的课程和有针对性的优化可能更加有效。随着文本到视频系统持续发展，PhyPrompt 这类提示词精炼方法为物理扎实的视频合成提供了一条参数高效的路径，可服务于机器人、科学可视化、教育等依赖物理真实性的应用。

% \subsection{Broader Impacts}
% \label{sec:impact}
% PhyPrompt offers clear positive benefits by enabling more physically plausible video synthesis, which can enhance applications in robotics, scientific visualization, virtual training, and education where realistic physics are crucial for downstream tasks. By improving semantic-and-physical fidelity, our met
